\documentclass[10pt,a4paper]{amsart}
\usepackage{amsmath,amssymb,amsthm,graphicx}
% Retained minimal preamble for the Stacks Project chapter more-algebra.tex (More on Algebra).
% The chapter file pinned-original-source.tex carries only \input{preamble} on its first line, so no
% preamble text lives inside the chapter itself.  The Stacks Project uses one single shared preamble for
% every chapter; the snapshot stages it identically as sources/stacks-chow/preamble.tex and
% sources/stacks-etale-cohomology/preamble.tex.  The block below is copied mechanically, line for line,
% from sources/stacks-chow/preamble.tex lines 14-18 (the Xy-pic setup) and 154-217 (the native theorem
% declarations together with the mathematical macros).  Nothing was added, removed or redefined; the
% wrapper class and the amsmath/amssymb/amsthm/graphicx base packages are supplied by the assembler.
% Chapter: more-algebra.tex -- More on Algebra.
% For commutative diagrams we use Xy-pic
\usepackage[all]{xy}

% We use 2cell for 2-commutative diagrams.
\xyoption{2cell}

% Theorem environments.
%
\theoremstyle{plain}
\newtheorem{theorem}[subsection]{Theorem}
\newtheorem{proposition}[subsection]{Proposition}
\newtheorem{lemma}[subsection]{Lemma}

\theoremstyle{definition}
\newtheorem{definition}[subsection]{Definition}
\newtheorem{example}[subsection]{Example}
\newtheorem{exercise}[subsection]{Exercise}
\newtheorem{situation}[subsection]{Situation}

\theoremstyle{remark}
\newtheorem{remark}[subsection]{Remark}
\newtheorem{remarks}[subsection]{Remarks}

\numberwithin{equation}{subsection}

% Macros
%
\def\lim{\mathop{\mathrm{lim}}\nolimits}
\def\colim{\mathop{\mathrm{colim}}\nolimits}
\def\Spec{\mathop{\mathrm{Spec}}}
\def\Hom{\mathop{\mathrm{Hom}}\nolimits}
\def\Ext{\mathop{\mathrm{Ext}}\nolimits}
\def\SheafHom{\mathop{\mathcal{H}\!\mathit{om}}\nolimits}
\def\SheafExt{\mathop{\mathcal{E}\!\mathit{xt}}\nolimits}
\def\Sch{\mathit{Sch}}
\def\Mor{\mathop{\mathrm{Mor}}\nolimits}
\def\Ob{\mathop{\mathrm{Ob}}\nolimits}
\def\Sh{\mathop{\mathit{Sh}}\nolimits}
\def\NL{\mathop{N\!L}\nolimits}
\def\CH{\mathop{\mathrm{CH}}\nolimits}
\def\proetale{{pro\text{-}\acute{e}tale}}
\def\etale{{\acute{e}tale}}
\def\QCoh{\mathit{QCoh}}
\def\Ker{\mathop{\mathrm{Ker}}}
\def\Im{\mathop{\mathrm{Im}}}
\def\Coker{\mathop{\mathrm{Coker}}}
\def\Coim{\mathop{\mathrm{Coim}}}

% Boxtimes
%
\DeclareMathSymbol{\boxtimes}{\mathbin}{AMSa}{"02}

%
% Macros for moduli stacks/spaces
%
\def\QCohstack{\mathcal{QC}\!\mathit{oh}}
\def\Cohstack{\mathcal{C}\!\mathit{oh}}
\def\Spacesstack{\mathcal{S}\!\mathit{paces}}
\def\Quotfunctor{\mathrm{Quot}}
\def\Hilbfunctor{\mathrm{Hilb}}
\def\Curvesstack{\mathcal{C}\!\mathit{urves}}
\def\Polarizedstack{\mathcal{P}\!\mathit{olarized}}
\def\Complexesstack{\mathcal{C}\!\mathit{omplexes}}
% \Pic is the operator that assigns to X its picard group, usage \Pic(X)
% \Picardstack_{X/B} denotes the Picard stack of X over B
% \Picardfunctor_{X/B} denotes the Picard functor of X over B
\def\Pic{\mathop{\mathrm{Pic}}\nolimits}
\def\Picardstack{\mathcal{P}\!\mathit{ic}}
\def\Picardfunctor{\mathrm{Pic}}
\def\Deformationcategory{\mathcal{D}\!\mathit{ef}}
\usepackage{cleveref}
\crefname{originallemma}{original lemma}{original lemmas}
\Crefname{originallemma}{Original lemma}{Original lemmas}
\makeatletter
\newlabel{lemma-fitting-ideals-and-pd1}{{0F7M}{source}{}{sourcepointer-1}{}}
\newlabel{lemma-fitting-ideals-and-pd1@cref}{{[originallemma][0][]0F7M}{[1][0][]source}{}{}{}}
\makeatother
\crefname{originalsection}{original section}{original sections}
\Crefname{originalsection}{Original section}{Original sections}
\makeatletter
\newlabel{section-eta}{{0F7N}{source}{}{sourcepointer-2}{}}
\newlabel{section-eta@cref}{{[originalsection][0][]0F7N}{[1][0][]source}{}{}{}}
\makeatother
\title{When the Fitting ideal of a bounded complex with respect to a nonzerodivisor is principal, the shifted complex is locally free and its differential splits off}
\author{The Stacks Project Authors}
\date{stacks-project:more-algebra.tex:local-snapshot-20261012; GNU FDL 1.2 or later}
\begin{document}
\maketitle
\setcounter{subsection}{1}
\noindent\textit{Editorial note.} All delivered mathematical text is the author's own, copied from the pinned Stacks Project chapter file more-algebra.tex; no statement or proof step was rewritten, reordered or removed. The counted statement is the labeled lemma environment at source lines 28049-28057 (canonical tag 0F7W) and the counted proof is the proof environment at source lines 28059-28102, which immediately follows it and is its own proof.\par\medskip
\noindent\textit{Editorial note.} The proof refers to the fitting ideal statement of this chapter and to the section in which the shifted complex is constructed; both are carried as navigation stubs with their roles recorded in the reasons.\par\medskip
\noindent\textit{Editorial note.} The statement uses the notation eta\_f for the shifted complex and the ideals I\_i(M, f) for the Fitting ideals, both as printed by the author in the surrounding section; the preamble supplies the macros involved.\par\medskip
\noindent\textit{Editorial note.} The chapter file pinned-original-source.tex carries only an input of the shared preamble on its first line, so no preamble text lives inside the chapter. The retained file sources/stacks-more-algebra/retained-preamble.tex is a mechanical, line-for-line copy of the project's single shared preamble as staged in the snapshot (Xy-pic setup plus the native theorem declarations and mathematical macros). Nothing was added, removed or redefined, and the excerpt serves only as preamble, never as a span coordinate.\par\medskip
\noindent\textit{Editorial note.} The source URL is the source-level Stacks Project URL, as required when more than one card is produced from one source; the per-claim tag is carried by claim\_key and by this note. The counted statement and proof cite no bibliography entry, so this card needs no bibliography.\par\medskip
\noindent\textit{Allowed standard prerequisites (editorial).}\begin{itemize}
\item The Fitting ideal I\_i(M, f) is generated by the r\_i by r\_i minors of the map (f, d\textasciicircum{}i) and its powers of f vanish after localizing at f. used at proof lines 28060-28068. Supplies the generator g and lets the author read the rank of the map after localizing, which gives the vanishing of the higher minors.
\item A finite module whose Fitting invariant is principal and which is defined by the minors of a map is finite locally free of the corresponding rank, and a finite projective module over a local ring is free. used at proof lines 28069-28078. Shows that the quotient Q is finite locally free of rank r\_\{i+1\} and f-torsion free, which identifies the cokernel modulo f-power torsion.
\item The shifted complex eta\_f M has as its i-th term the kernel of the map (d\textasciicircum{}i, -1) from f\textasciicircum{}i M\textasciicircum{}i direct sum f\textasciicircum{}\{i+1\} M\textasciicircum{}\{i+1\} to f\textasciicircum{}i M\textasciicircum{}\{i+1\}. used at proof lines 28094-28098. Identifies the term of the shifted complex with the kernel in question and thereby with a direct summand of the finite free module.
\end{itemize}\medskip
\ifdefined\hypertarget\hypertarget{sourcepointer-1}{}\fi
\noindent\textit{Source reference.} Original lemma 0F7M: sources/stacks-more-algebra/pinned-original-source.tex, source lines 1616--1616. The reference is applied at proof lines 28075-28076 to conclude that the quotient by g-torsion is finite locally free of rank r\_\{i+1\}; the identification of the shifted term as a direct summand, which is the conclusion, is obtained inside the counted span.\par\medskip
\ifdefined\hypertarget\hypertarget{sourcepointer-2}{}\fi
\noindent\textit{Source reference.} Original section 0F7N: sources/stacks-more-algebra/pinned-original-source.tex, source lines 27385--27385. The reference points to the construction of the shifted complex, used at proof lines 28093-28096 to identify the image of the injective map with the kernel of the second differential; it is a pointer to the definition, not a substitute for the computation in the span.\par\medskip
\begin{lemma}
\label{lemma-eta-locally-free}
Let $A$ be a ring and let $f \in A$ be a nonzerodivisor. Let
$M^\bullet$ be a bounded complex of finite free $A$-modules. Assume
$I_i(M^\bullet, f)$ is a principal ideal.
Then $(\eta_fM)^i$ is locally free of rank $r_i$ and the map
$(1, d^i) : (\eta_fM)^i \to f^iM^i \oplus f^{i + 1}M^{i + 1}$
is the inclusion of a direct summand.
\end{lemma}

\begin{proof}
Choose a generator $g$ for $I_i(M^\bullet, f)$. Since
$f^{r_i} \in I_i(M^\bullet, f)$ we see that $g$ divides a power of $f$.
In particular $g$ is a nonzerodivisor in $A$.
The $r_i \times r_i$-minors of the map
$(f, d^i) : M^i \to M^i \oplus M^{i + 1}$
generate the ideal $I_i(M^\bullet, f)$ and the
$(r_i + 1) \times (r_i + 1)$-minors of $(f, d^i)$ are zero:
we may check this after localizing at $f$
where the rank of the map is equal to $r_i$.
Consider the surjection
$$
M^i \oplus M^{i + 1}
\longrightarrow
Q = \Coker(f, d^i)/g\text{-torsion}
$$
By Lemma \ref{lemma-fitting-ideals-and-pd1}
the module $Q$ is finite locally free of rank $r_{i + 1}$.
Hence $Q$ is $f$-torsion free and we conclude the cokernel of
$(f, d^i)$ modulo $f$-power torsion is $Q$ as well.

\medskip\noindent
Consider the complex of finite free $A$-modules
$$
0 \to f^{i + 1}M^i \xrightarrow{1, d^i} f^iM^i \oplus f^{i + 1}M^{i + 1}
\xrightarrow{d^i, -1} f^iM^{i + 1} \to 0
$$
which becomes split exact after localizing at $f$.
The map $(1, d^i) : f^{i + 1}M^i \to f^iM^i \oplus f^{i + 1}M^{i + 1}$
is isomorphic to the map $(f, d^i) : M^i \to M^i \oplus M^{i + 1}$ we
studied above. Hence the image
$$
Q' = \Im(f^iM^i \oplus f^{i + 1}M^{i + 1} \xrightarrow{d^i, -1} f^iM^{i + 1})
$$
is isomorphic to $Q$ in particular projective. On the other hand,
by construction of $\eta_f$ in Section \ref{section-eta}
the image of the injective map
$(1, d^i) : (\eta_fM)^i \to f^iM^i \oplus f^{i + 1}M^{i + 1}$ is
the kernel of $(d^i, -1)$. We conclude that we obtain an
isomorphism $(\eta_fM)^i \oplus Q' = f^iM^i \oplus f^{i + 1}M^{i + 1}$
and we see that indeed $\eta_fM^i$ is finite locally free
of rank $r_i$ and that $(1, d^i)$ is the inclusion of a direct
summand.
\end{proof}

\end{document}
